% =====================================================================
%  Chapter 3 — Methodology and System Design  (~2,500 words)
%  DRAFT generated 2026-07-18. Re-threaded from week2_gating.tex — sprint
%  framing removed, equations preserved. Needs: architecture figure
%  (fig placeholder marked below).
% =====================================================================
\chapter{Methodology and System Design}
\label{ch:method}

\section{Design Overview}
\label{sec:method-overview}

The system is organised as a six-layer pipeline covering configuration, data
contracts, ingestion, models and retrieval, policies and gating, and
evaluation, and one principle governs the whole structure, which is that each
layer depends only on the layers beneath it and communicates exclusively
through shared typed dataclasses (Figure~\ref{fig:architecture}). Every
component can therefore be unit-tested in isolation against mock
collaborators, and the retrieval gates were developed against a stable
test-covered substrate rather than a moving target. The same discipline is
what makes the evaluation trustworthy, because every policy emits an
identical fully populated record per query, so analyses such as the paired
retrieval-distraction study can be run entirely from logs.

\begin{figure}[t]
  \centering
  \begin{tikzpicture}[
    font=\small,
    layer/.style={draw, rounded corners=2pt, minimum width=10.4cm,
                  minimum height=0.82cm, align=center, fill=black!3},
    novel/.style={layer, fill=blue!8, very thick},
    arrow/.style={-{Latex[length=2mm]}, thick, black!55},
  ]
    \node[layer] (cfg)  at (0, 0)     {\textbf{Configuration} - layered YAML $\rightarrow$ validated \texttt{ProjectConfig}};
    \node[layer] (data) at (0,-1.05)  {\textbf{Data contracts} - \texttt{UnifiedQuestion}, \texttt{PolicyResult}, \texttt{RunRecord}};
    \node[layer] (ing)  at (0,-2.10)  {\textbf{Ingestion} - dataset loaders, chunking, risk tagger};
    \node[layer] (mod)  at (0,-3.15)  {\textbf{Models \& retrieval} - \texttt{llama.cpp} backend $\mid$ BM25, dense, hybrid (RRF)};
    \node[novel] (gate) at (0,-4.20)  {\textbf{Policies \& gating} - P1--P5 $\mid$ entropy, margin, probe $\rightarrow$ majority vote};
    \node[layer] (eval) at (0,-5.25)  {\textbf{Evaluation} - scorers, profiler, append-only JSONL audit log};

    \draw[arrow] (cfg)  -- (data);
    \draw[arrow] (data) -- (ing);
    \draw[arrow] (ing)  -- (mod);
    \draw[arrow] (mod)  -- (gate);
    \draw[arrow] (gate) -- (eval);

    \draw[arrow, dashed] (eval.east) -- (5.9,-5.25) -- (5.9,-4.20) -- (gate.east);
  \end{tikzpicture}
  \caption{The six-layer architecture. Each layer depends only on those
    above it and communicates through the shared typed dataclasses, so every
    component is unit-testable against mock collaborators. The gating layer
    (highlighted) is the novel contribution. Because the evaluation layer
    logs each query's full gate payload, the analyses of
    Chapter~\ref{ch:eval}, namely threshold calibration, gate ablations and
    the paired distraction study, replay from those logs without re-invoking
    the model (dashed path).}
  \label{fig:architecture}
\end{figure}

\section{The Policy Space}
\label{sec:method-policies}

A \emph{policy} is a complete answering strategy: given a question, it decides
what (if anything) to retrieve and produces an answer. All policies implement
a single method, \texttt{answer(question) $\rightarrow$ PolicyResult}, so the
run driver is agnostic to which policy it holds, which is the precondition for a
fair comparative experiment. Four policies are evaluated
(Table~\ref{tab:setup}):

% AUTO-GENERATED by scripts/generate_tables.py -- DO NOT EDIT.
% Regenerate with: python scripts/generate_tables.py
% Every number is computed from results/raw_logs/ via evaluation.loading.
% git b3a9d01 | 2026-07-13 13:27
\begin{table}[t]
  \centering
  \begin{tabular}{llll}
    \toprule
    Policy & Retrieval & Gate & Role in the evaluation \\
    \midrule
    P1 always-retrieve & every query & --- & retrieval baseline \\
    P4 hybrid & every query & --- & BM25\,+\,dense fusion (RRF) \\
    P3 closed-book$^{\dagger}$ & never & --- & parametric-knowledge reference \\
    P5 gated & selective & 3-gate ensemble & the proposed method \\
    \bottomrule
  \end{tabular}
  \caption{The four policies compared. P5 decides per query whether to retrieve,
  using a majority vote ($\geq$2 of 3) over the token-entropy, logit-margin and
  hallucination-probe gates. $^{\dagger}$P3 never retrieves and is therefore
  corpus-independent.}
  \label{tab:setup}
\end{table}


\begin{itemize}
  \item \textbf{P1 (always-retrieve)} retrieves for every query and answers
    with the evidence in context. It is the cost ceiling and the conventional
    RAG baseline.
  \item \textbf{P3 (closed-book)} never retrieves: the parametric-knowledge
    reference point and the cost floor.
  \item \textbf{P4 (hybrid)} is P1 bound to the hybrid (rank-fusion)
    retriever rather than BM25 alone.
  \item \textbf{P5 (gated)}, the proposed method, consults the gate
    ensemble per query and behaves like P1 on \textsc{retrieve} and like P3 on
    \textsc{skip}.
\end{itemize}

A fifth policy, \textbf{P2 (always-retrieve with citations)}, is implemented,
extending P1 by asking the model to name its sources, matched leniently back to
retrieved chunks after the attribution-evaluation convention of
\citet{alce2023}, but is not evaluated as a distinct condition here, for two
reasons. Its retrieval behaviour is identical to P1, so its answer accuracy
carries no information the central ``when to retrieve'' question does not already
get from P1. And its one distinct output, citation quality, cannot be scored on
these benchmarks: MIRAGE supplies only a gold letter and PubMedQA a free-text
paragraph, neither of which provides the passage-level attribution gold that
grounded-citation precision and recall require. The source-transparency goal P2
serves, making the evidence behind an answer inspectable is nonetheless
realised qualitatively in the demo interface (Section~\ref{sec:impl-demo}),
which surfaces each retrieved chunk with the query terms it matched; what P2 as
a \emph{scored} policy would add is a quantitative citation precision and recall,
and that number, requiring an ALCE-style annotated corpus, is left to future
work (Section~\ref{sec:conc-future}). Two further policies from the original plan
were scoped out entirely: a user-facing retrieval toggle (P6) and a multi-agent
variant (P7). Both add breadth without informing when a system \emph{should}
retrieve, and the evaluation depth spent instead on paired analyses and
ablations (Chapter~\ref{ch:eval}) proved the better trade.

\section{Retrieval}
\label{sec:method-retrieval}

Three retrievers are implemented behind a common interface, selected by
configuration rather than code change.

\paragraph{Lexical (BM25).} BM25 \citep{robertson2009bm25} scores documents
by weighted lexical overlap. It requires no embedding model, no vector index
and no GPU, making it the only retriever viable on the lowest hardware tier, and
medical terminology is specific enough that exact lexical match is a strong
signal (drug names, dosages, anatomical terms).

\paragraph{Dense (vector).} Each corpus chunk is embedded once with
\texttt{all-MiniLM-L6-v2} \citep{reimers2019} (384 dimensions) and indexed
in a FAISS \texttt{IndexFlatIP} \citep{johnson2019faiss}. Embeddings are
L2-normalised, so inner product equals cosine similarity; the index is exact
(brute-force), so retrieval is deterministic and needs no training, which is a
reproducibility choice made deliberately, accepting linear-scan cost in
exchange for run-to-run identical rankings. Dense retrieval catches the
paraphrase cases BM25 misses: a query about a \emph{heart attack} matching a
passage that only says \emph{myocardial infarction}.

\paragraph{Hybrid (RRF).} The two retrievers have complementary strengths,
and their scores live on incomparable scales. Reciprocal Rank Fusion merges
them scale-free by fusing \emph{ranks}:
\begin{equation}
  \mathrm{score}(d) \;=\; \sum_{i \in \{\text{bm25},\,\text{vec}\}}
  \frac{1}{k + \mathrm{rank}_i(d)}, \qquad k = 60,
  \label{eq:rrf}
\end{equation}
where $\mathrm{rank}_i(d)$ is the 1-based rank of document $d$ under
retriever $i$. Documents ranked highly by both rise to the top; no score
normalisation is ever required. The damping constant $k=60$ is the standard
value and is exposed in configuration.

\section{The Gate Ensemble}
\label{sec:method-gates}

\subsection{Why Three Signals}
\label{sec:method-why-three}

A single uncertainty signal is fragile, because each of the signal families
surveyed in Section~\ref{sec:bg-calibration} carries a characteristic
failure mode under which it would gate incorrectly.

\emph{Token entropy} is inflated by benign lexical freedom, because
probability mass spreads wherever multiple phrasings are equally acceptable,
such as formatting tokens, function words and punctuation, regardless of
whether the underlying fact is certain. A model that knows the answer but has
many ways to phrase it therefore looks uncertain to an entropy gate, and the
token-level attribution of Section~\ref{sec:eval-attribution} demonstrates
exactly this failure on real clinical text.

\emph{Logit margin} fails in the opposite direction, since its failure mode is
confident error. A model that has internalised a wrong fact produces a large
margin on a wrong answer and the gate reads certainty, because margin measures
decisiveness rather than correctness.

\emph{Verbalized confidence}, meaning asking the model to rate its own answer,
is the standard third option, but \citet{kadavath2022} show that
self-assessment calibration degrades with model size, so self-reports from a
3B model are close to uninformative. This project therefore replaces it with a
stability test in the spirit of self-consistency
\citep{wang2023selfconsistency}, generating two drafts under different prompt
framings and measuring their agreement, on the reasoning that genuinely held
knowledge is invariant to surface framing while a guess is not. Reduced from
many-sample majority voting to a two-draft agreement test, the idea becomes
cheap enough to run per query on CPU.

The design response to these three failure modes is an ensemble drawing on
different aspects of the model's behaviour, namely distributional spread,
decision-boundary decisiveness and cross-framing stability, combined by
majority vote so that no single failure mode dominates. The framing is
safety-asymmetric by intent, since an unnecessary retrieval costs seconds of
latency while a confidently wrong closed-book answer in a clinical setting
costs far more, so the ensemble is deliberately biased toward retrieving when
its signals disagree. The evaluation partially complicates this story, because
the ablation of Section~\ref{sec:eval-ablation} finds that the probe
over-triggers retrieval and that the flat majority vote is what protects the
ensemble from it, and that result is reported there rather than smoothed over
here.

\subsection{Gate 1: Token Entropy}
\label{sec:method-gate-entropy}

The entropy gate generates a short draft answer (48 tokens) without retrieval
and measures uncertainty from the full-vocabulary distribution at each
generated position. For generated token $t$ with softmax distribution $p_t$
over vocabulary $V$, the per-token Shannon entropy (nats) and gate signal are
\begin{align}
  H(p_t) &= -\sum_{v \in V} p_t(v)\,\log p_t(v), &
  \bar{H} &= \frac{1}{N}\sum_{t=1}^{N} H(p_t),
  \label{eq:entropy}
\end{align}
and the gate votes \textsc{retrieve} when $\bar{H} > \tau_H$. High mean
entropy indicates probability mass spread over many continuations, that is, direct
lexical-level uncertainty.

Beyond the scalar decision signal, the gate returns the full per-token array
$\{H(p_t)\}_{t=1}^N$ and propagates it into the audit log. This array is what
enables the token-level entropy attribution analysis
(Section~\ref{sec:eval-attribution}): uncertainty can be localised to
individual tokens of the draft, not merely summarised. The gate requires the
raw logit buffer (\texttt{logits\_all=true}); where that buffer is
unavailable it \emph{abstains} rather than fabricating a signal.

\subsection{Gate 2: Logit Margin}
\label{sec:method-gate-margin}

The margin gate measures decisiveness from the gap between the two most
probable tokens at each position. For top-two probabilities
$p_t^{(1)} \ge p_t^{(2)}$,
\begin{align}
  m_t &= p_t^{(1)} - p_t^{(2)}, &
  \bar{M} &= \frac{1}{N}\sum_{t=1}^{N} m_t,
  \label{eq:margin}
\end{align}
and the gate votes \textsc{retrieve} when $\bar{M} < \tau_M$. A small margin
means the model is torn between near-equiprobable continuations, making it
complementary to entropy, which summarises the whole distribution while the
margin isolates the decision boundary. Practically, the margin gate reads the
completion API's top-$k$ log-probabilities rather than the raw logit buffer:
a lighter dependency that survives \texttt{logits\_all=false}, and one whose
parser was validated against the real inference library so that the unit-test
mock is faithful to production output.

\subsection{Gate 3: Hallucination Probe}
\label{sec:method-gate-probe}

The probe replaces the verbalized-confidence gate of the original design.
Rather than asking a small model to introspect, a task at which it is
unreliable, it tests whether the model's answer is \emph{stable under
reframing}. Two drafts are generated from the same question under two prompt
framings:
\begin{align*}
  \text{draft}_A &= \mathcal{M}\bigl(\text{probe}_A(q)\bigr), &
  \text{draft}_B &= \mathcal{M}\bigl(\text{probe}_B(q)\bigr),
\end{align*}
where $\text{probe}_A$ uses an assistant-role framing and $\text{probe}_B$ a
direct-question framing. For multiple-choice questions the gate compares
extracted option letters and votes \textsc{retrieve} when they differ; for
open-ended questions it computes token-level $F_1$ between the drafts and
votes \textsc{retrieve} when $F_1 < 0.7$. The intuition: genuinely held
knowledge is invariant to surface framing; a guess is sensitive to it.

The probe is purely text-based, since it never touches the logit buffer, and
is therefore the only gate available on the lowest hardware tier. Its cost is
two extra draft generations per query, which turns out to be the dominant
term in the gated policy's latency (Section~\ref{sec:eval-cost}); that cost
is reported honestly rather than amortised away.

\section{Ensemble Voting and Hardware-Aware Degradation}
\label{sec:method-ensemble}

The gates are combined by majority vote over \emph{available} members. Let
$g_i \in \{\textsc{retrieve}, \textsc{skip}\}$ be the decision of member $i$
among the available set
$A \subseteq \{\text{entropy}, \text{margin}, \text{probe}\}$. The ensemble
retrieves when
\begin{equation}
  \bigl|\{\, i \in A : g_i = \textsc{retrieve} \,\}\bigr| \;\ge\; v_{\min},
  \qquad v_{\min} = 2 .
  \label{eq:vote}
\end{equation}
Only members reporting themselves available cast votes; an abstaining gate is
excluded from both the count of retrieve votes and the electorate.

This availability accounting yields a principled \emph{degraded mode}. On the
low hardware tier only the probe is available, so $v_{\min}=2$ is
unreachable; the ensemble then falls back to ``retrieve if any available
member votes retrieve'' and records a \texttt{degraded} flag in the audit
log. The fallback is enforced at \emph{construction} time: the policy factory
inspects the model's \texttt{logits\_all} setting and downgrades any
requested gate to probe-only when the logit buffer is disabled. The gate a
policy actually runs is therefore constrained by hardware, not merely by
configuration, which is an auditable property.

An honest scope statement belongs here: the degradation \emph{mechanism} is
implemented and unit-tested, but only the medium hardware tier was physically
benchmarked in this study. Per-tier latency and energy on real 4\,GB and
16\,GB devices remain future work
(Section~\ref{sec:disc-limitations}).

\section{The Safety Envelope}
\label{sec:method-envelope}

The evaluation's organising concept for clinical acceptability is the
\emph{safety envelope}. Each question carries a rule-assigned clinical risk
level (low, medium, high; Section~\ref{sec:impl-config}); each risk level
carries a stipulated accuracy bar (\barLow{}, \barMedium{}, \barHigh{}
respectively). For each risk level, the envelope asks: \emph{which is the
cheapest policy whose accuracy clears the bar?} A policy inside the envelope
is, by this operationalisation, cheap enough and safe enough to consider for
that stratum.

The bars are stipulated, not derived from clinical practice, because there is no
consensus mapping from benchmark accuracy to clinical safety, and this
dissertation does not pretend to one. The contribution is the \emph{method}
of asking the question at all, plus the honest answer it returns at this
model scale: the envelope is empty (Section~\ref{sec:eval-envelope}), and its
emptiness is a measurement of the distance still to travel, not a failure to
be hidden.